\section{Resolved difficulties, retained hypotheses, and limits}\label{sec:limits}

\subsection{What the framework resolves}

\begin{longtable}{@{}P{.25\textwidth}P{.35\textwidth}P{.32\textwidth}@{}}
\toprule Difficulty & Positive result & Data still required\\\midrule\endhead
Finite solutions at each scale need not form a point. & Residual diameter control gives a Cauchy name, canonical residual value zero, and uniqueness in the completion. & An actual PR approximate-solution routine and a proved PR diameter modulus.\\
A Cauchy estimate does not identify the value of a function at the limit. & Quantitative extension, density, and uniqueness identify the completed evaluation. & A proved modulus and extensional finite evaluator.\\
Local chart selection can be discontinuous or undecidable. & Weighted gluing is continuous with constant \(L+2NM/\delta\); safe inner weights handle named distances. & A positive denominator certificate and local correctness on every used support.\\
Finite sample tables may not be realizable as legitimate functions. & Every clipped envelope is a bounded Lipschitz function; quantized tables give an explicit net. & Shared range/continuity bounds; evaluator data when selecting a center near an input function.\\
Quantification ranges over continuous objects and functions. & Total correspondences yield finite continuous-formula values with no quantifier-depth error loss for the core connectives. & Compact presentations, atomic moduli, and finite evaluators.\\
Derivative separation fails at a singular inverse. & A H\"older residual modulus constructs the real cube root through zero. & Specific cubic separation and approximate-solvability estimates.\\
Continuous partitions of unity do not preserve holomorphicity. & Exact overlap equality makes the glued function locally equal to a certified analytic chart. & Power-series majorants, positive chart/domain margins, and all-precision overlap equality.\\
An infinite-domain optimization problem seems to demand choice. & The tent problem has exact value \(1/4\), an explicit attainer, and residual modulus \(2k+3\). & A concrete retraction and quantitative uniqueness proof; these are supplied for this example.\\
\bottomrule
\end{longtable}

\subsection{What is not implied}

The results do not turn every true theorem of analysis into an \ipra\ theorem. Qualitative continuity does not supply a PR modulus. A true existential statement does not supply a PR witness or a PR search bound. The logical complexity of a finite verification predicate must be addressed by actual bounded Boolean compilation, and the loop invariant must be proved with quantifier-free induction. These are mathematical obligations, not implementation details to omit.

Neither total boundedness nor finite solvability at every scale selects a computable infinite branch in general. Our function-space construction supplies finite nets and a completion operation for already modulated Cauchy sequences; it does not extract subsequences from arbitrary sequences. Computing a supremum value is distinct from computing a point attaining it. The residual theorem adds exactly the quantitative uniqueness information used to bridge that distinction in the stated examples.

The model-theoretic results evaluate a fixed formula on a supplied structure or compare supplied structures. They do not construct a model of an arbitrary consistent theory, prove first-order compactness inside \ipra, supply ultraproducts, or decide exact zero-valued truth. The observable-quotient example works only for a language whose atoms factor through the specified projection; adding the full supremum norm of the original noncompact space invalidates that premise.

The categorical results are equally scoped. Completion has a universal extension property in the external category of the chosen presentations. Finite correspondences compose with an additive error. Fixed Lipschitz bounds support the stated evaluation/currying correspondence. These facts do not assert a category of all sets or all structures inside first-order arithmetic, unrestricted exponentials, arbitrary quotients, or arbitrary colimits. A setoid interpretation needs derivations of the translated source axioms, as the existing ProveIt interface already requires.

Finally, the successful singular real cube root does not remove the known obstruction to uniform branch selection. The complex cube root has no single holomorphic inverse around zero, and monodromy can prevent global analytic gluing even when every chart is valid. Our analytic descent theorem retains the exact overlap premise rather than inferring it from a finite cover.

\subsection{A practical interface for future packets}

An analytic or abstract packet fits this framework when it explicitly supplies five components:
\begin{enumerate}
\item A numerical presentation of the relevant objects and finite observations, including domain and denominator guards.
\item Fixed PR evaluators or finite-table algorithms, with moduli, loop bounds, and nonempty finite supports where needed.
\item Finite correctness and stability inequalities, with quantifier-free or Booleanized bounded invariants.
\item A coherence or representation-independence argument, including any diameter, separation, or overlap modulus.
\item An adequacy theorem identifying the resulting names or formula values with the desired continuous meaning.
\end{enumerate}

This checklist does not add a new axiom. It locates the mathematical work and makes reuse possible. For the present report the most valuable next implementation is the ordinary-iPRA finite correspondence theorem, followed by the bridge from the existing coherent-minimizer name to its canonical completed residual. The compiled finite kernel and the explicit numerical schedules make those next steps concrete without claiming that the remaining deep formalization has already been completed.
